% SPDX-FileCopyrightText: 2026 Will Cook
% SPDX-License-Identifier: CC-BY-4.0
%
% One of the Erdős Problem Notes: a short standalone treatment of a single
% problem in the ErdosProblems expansion library.  Build with
% `tectonic erdos-1049-rational-base-lambert.tex` or `make`.
\documentclass[11pt]{article}
% SPDX-FileCopyrightText: 2026 Will Cook
% SPDX-License-Identifier: CC-BY-4.0
%
% Shared preamble for the Erdős Problem Notes: the short, standalone,
% problem-owned manuscripts covering the expansion library `ErdosProblems`.
%
% One note owns one Erdős problem.  Each note repeats whatever context its own
% reader needs, so that a reader who arrives at a single problem never has to
% assemble it from the gateway paper.  Deliberate repetition across notes is the
% design, not an oversight.
%
% Source links resolve against one pinned formal-source commit, declared once
% below.  `scripts/check_problem_note_sources.py` verifies that every authored
% (file, line, declaration) triple names that exact declaration in the pinned
% snapshot, so a link cannot silently rot as later work moves lines.

\usepackage{paper-house-style}
\usepackage{array}
\usepackage{booktabs}
\usepackage{enumitem}
\setlist{topsep=0.35em,itemsep=0.2em}
\usepackage[colorlinks=true,linkcolor=linkinternal,urlcolor=linkexternal,
  citecolor=linkinternal]{hyperref}
\hypersetup{bookmarksnumbered=true,bookmarksopen=true,bookmarksopenlevel=1,
  pdfauthor={Will Cook},pdflang={en-GB}}

% ---------------------------------------------------------------------------
% Pinned formal source.  \commit names the checkpoint every link resolves
% against; the checker reads that snapshot rather than the working tree, so the
% printed coordinates stay correct however later waves move declarations.
% ---------------------------------------------------------------------------
\newcommand{\commit}{e5fd26a6d1b1a346430205eab5385b4c9565d004}
\newcommand{\commitshort}{e5fd26a6d1b1}
\newcommand{\repobase}{https://github.com/wcook04/plectis-lean-erdos249-257/blob/\commit}
\newcommand{\sourceurl}{https://github.com/wcook04/plectis-lean-erdos249-257/tree/\commit}
\newcommand{\PX}{ErdosProblems}

% Declaration names stay inside link targets.  The printed label is either a
% spaced, lower-case reading of the declaration name (\lref) or a short authored
% phrase (\lword); neither prints a module path or a raw Lean identifier.
\ExplSyntaxOn
\cs_new_protected:Npn \noteleanlabel #1
  {
    \tl_set:Nx \l_tmpa_tl { \tl_to_str:n {#1} }
    \regex_replace_all:nnN { _+ } { \c{space} } \l_tmpa_tl
    \regex_replace_all:nnN
      { ([a-z0-9])([A-Z]) }
      { \1\c{space}\2 }
      \l_tmpa_tl
    \tl_set:Nx \l_tmpa_tl { \tl_lower_case:n { \tl_use:N \l_tmpa_tl } }
    \tl_use:N \l_tmpa_tl
  }
\ExplSyntaxOff
\newcommand{\leanlabel}[1]{\noteleanlabel{#1}}
\newcommand{\lref}[3]{\href{\repobase/\PX/#1\#L#2}{\leanlabel{#3}}}
\newcommand{\lword}[4]{\href{\repobase/\PX/#1\#L#2}{#4}}
\newcommand{\lloc}[2]{\href{\repobase/\PX/#1\#L#2}{Lean source}}

% Links into a sibling library, named repository-relative.  The expansion
% library is implicit in \lref/\lword, because that is where problem-owned
% work lands.  Several problems have their headline theorem in the older
% reviewed #249/#257 corpus, and a note that could not name that library
% would have to paraphrase a checked statement instead of linking it.
\newcommand{\mref}[3]{\href{\repobase/#1\#L#2}{\leanlabel{#3}}}
\newcommand{\mword}[4]{\href{\repobase/#1\#L#2}{#4}}
\newcommand{\mloc}[2]{\href{\repobase/#1\#L#2}{Lean source}}
\emergencystretch=2em

\newtheorem{theorem}{Theorem}[section]
\newtheorem{proposition}[theorem]{Proposition}
\newtheorem{lemma}[theorem]{Lemma}
\newtheorem{corollary}[theorem]{Corollary}
\theoremstyle{definition}
\newtheorem{definition}[theorem]{Definition}
\newtheorem{example}[theorem]{Example}
\newtheorem{problem}[theorem]{Problem}
\theoremstyle{remark}
\newtheorem*{remark}{Remark}

\newcommand{\N}{\mathbb{N}}
\newcommand{\Npos}{\mathbb{N}_{>0}}
\newcommand{\Z}{\mathbb{Z}}
\newcommand{\Q}{\mathbb{Q}}
\newcommand{\R}{\mathbb{R}}
\newcommand{\lcm}{\operatorname{lcm}}
\newcommand{\ph}{\varphi}

% The series line printed under every note's title block.
\newcommand{\seriesline}{Erd\H{o}s Problem Notes}

% Production provenance.  The wording is fixed by the corpus provenance
% contract and reproduced verbatim in every manuscript in this repository, so
% the papers cannot drift into different accounts of how they were made.
\newcommand{\noteprovenance}{\provenancenote{\emph{Authorship and AI use.} The
prose, formal proofs, and software in this version were drafted and revised by
large-language-model agents under Will Cook's direction.  Cook set the
objectives and acceptance criteria, reviewed and authorised the manuscript, and
is responsible for its contents.  The agents are production tools, not authors.
See \emph{Statements and declarations}.}}

% The status paragraph every note carries, in its own words, before any result.
% Kernel checking establishes the exact Lean proposition; it does not establish
% that the proposition is the interesting one, that it is new, or that the
% Erdős problem is closed.
\newcommand{\statusboundary}{%
  \emph{Status.}  The problem treated here is open, and this note does not
  close it.  Every statement below marked as checked is a proposition that the
  pinned Lean kernel accepts from the sources this note links to, with no
  \texttt{sorry}, no added axiom, and no unchecked evaluation.  That is a claim
  about the formal statement, not about its mathematical interest, its
  novelty, or the original problem.  The unresolved obligations are named
  exactly, in their own section, and none of the finite computations,
  reductions, or no-go results here removes one of them.}

\hypersetup{
  pdftitle={Rational bases and the integer-base argument: a Lean-checked strategy separation for Erdős Problem 1049},
  pdfsubject={Erdős Problem Notes: Problem 1049},
  pdfkeywords={irrationality, Lambert series, rational base, Padé approximation, Lean 4},
}

\runninghead{Erd\H{o}s \#1049: rational bases}
\title{Rational Bases and the Integer-Base Argument}
\subtitle{A Lean-checked strategy separation for Erd\H{o}s Problem~\#1049}
\author{Will Cook}
\date{July 2026}

\begin{document}
\maketitle
\noteprovenance

\begin{abstract}
\noindent
Erd\H{o}s's argument for the irrationality of $\sum_{n\ge1}c(n)b^{-n}$ at an
integer base $b\ge2$ proceeds by clearing denominators at a cut level and
trapping a positive integer below $1$.  At a rational base $r/s$ with $s\ge2$
the same clearing leaves a factor $s^{N+1}$ behind.  We formalise in Lean the
elementary consequence of that residue.  Writing a coordinatewise corridor for
the surviving arithmetic, we prove that a corridor forces
$s^{N+K+1}<r(N+K)$, and that at the base $3/2$ this is impossible once the
shift and the cleared window are both nonempty.  We also prove the exact
denominator-cleared tail recurrence
$U_{N+1}=rU_N-Bc(N+1)s^{N+1}$ and the bound
$2^{N+1}\le Bc(N+1)s^{N+1}$ for $s\ge2$ and positive data, against the exact
collapse of that term to $Bc(N+1)$ at $s=1$.

Problem~\#1049 is open, and nothing here decides it.  The results are
statements about one clearing scheme, not about the arithmetic nature of any
value: we do not prove irrationality at $3/2$, at any other rational base, or
in the unrestricted problem.  The published height criterion that does settle
a restricted family of rational bases, including $7/2$, is used as an external
theorem; of its inputs we formalise only two exact arithmetic facts, and we
state plainly which three inputs are not formalised.  Separately we check
polynomial bounds on the denominator exponents of a homogenised Pad\'e
construction, with the analytic remainder left untouched.
\end{abstract}

\section{The problem}\label{sec:problem}

Let $t>1$ be a rational number and let $\tau(n)$ count the divisors of $n$.
Erd\H{o}s Problem~\#1049 asks whether
\[
 F(t)=\sum_{n\ge1}\frac{1}{t^{n}-1}=\sum_{n\ge1}\frac{\tau(n)}{t^{n}}
\]
is irrational~\cite[p.~102]{erdos1988}.  The two forms agree by expanding
$(t^{n}-1)^{-1}=\sum_{k\ge1}t^{-nk}$ and collecting the terms with the same
exponent, the coefficient of $t^{-m}$ being the number of divisors of $m$.
The question is a conjecture of Chowla; Erd\H{o}s proved it for every integer
$t\ge2$~\cite{erdos1948}.  Numbering and status follow Bloom's
catalogue~\cite{erdosproblems}.  For non-integer rational $t$ the problem is
open.

Write $t=r/s$ in lowest terms with $r>s\ge1$, so that $s=1$ is exactly the
integer case Erd\H{o}s settled.  The smallest resistant explicit base is
$t=3/2$.  A published height criterion of Bundschuh and
V\"a\"an\"anen~\cite{bv1994} settles a family of rational bases restricted by a
height condition; that family contains $7/2$ and does not contain $3/2$.
Between the two lies the question this note is about: what exactly stops the
integer-base argument from running at $3/2$?

\paragraph{Relation to prior work.}
The statement of Problem~\#1049 has been formalised before, as a conjecture
with an unfilled proof, in the \emph{formal-conjectures} collection.  That is a
statement of the question; the present note formalises propositions about the
clearing argument, and neither answers it.  Prior mathematical work on the
problem is Erd\H{o}s's integer-base theorem~\cite{erdos1948} and the height
criterion of~\cite{bv1994}; no claim of priority is made for anything below,
which concerns the formal status of an argument rather than a new theorem about
$F(t)$.

\paragraph{The integer-base argument, in outline.}
Suppose $F(b)=p/q$.  Cut the series at a level $N$, multiply by $q b^{N}$, and
subtract the cleared prefix.  What remains is a positive integer, since every
term of the prefix has been cleared and the tail is positive; and it is smaller
than $1$, since the tail of a geometric-type series is small.  A positive
integer below $1$ does not exist.  The argument is elementary and it is
sensitive in exactly one place: the clearing must remove every denominator.

At $\beta=r/s$ the term $c(n)\beta^{-n}$ is $c(n)s^{n}/r^{n}$.  Clearing the
power of $r$ leaves the numerator factor $s^{n}$ in place.  That factor is
invisible when $s=1$ and grows geometrically when $s\ge2$.  The two sections
that follow record what that residue costs, first as a no-go for the clearing
scheme (Section~\ref{sec:corridor}) and then as an exact recurrence with a
lower bound on the surviving term (Section~\ref{sec:tail}).

\paragraph{Status of everything below.}
\statusboundary

\begin{table}[t]
\centering
\small
\begin{tabular}{@{}
  >{\raggedright\arraybackslash}p{0.34\linewidth}
  >{\raggedright\arraybackslash}p{0.18\linewidth}
  >{\raggedright\arraybackslash}p{0.40\linewidth}@{}}
\toprule
Statement & Status & Treatment here \\
\midrule
Irrationality of $F(3/2)$ & Open & Not proved, and no partial result here bears on it. \\
Irrationality at some rational non-integer base & Proved elsewhere & The height criterion of~\cite{bv1994}; cited, not formalised. \\
Corridor forces $s^{N+K+1}<r(N+K)$ & Proved here & Theorem~\ref{res:corridorbound}. \\
No corridor at base $3/2$ & Proved here & Theorem~\ref{res:nocorridor}; a statement about the clearing scheme. \\
Cleared-tail recurrence & Proved here & Theorem~\ref{res:tailrec}, an exact identity. \\
Forcing term at least $2^{N+1}$ for $s\ge2$ & Proved here & Theorem~\ref{res:forcing}; at $s=1$ the term is $Bc(N+1)$. \\
Two arithmetic facts behind the $7/2$ criterion & Proved here & Section~\ref{sec:sevenhalves}; three further inputs are not formalised. \\
Pad\'e denominator-exponent bounds & Proved here & Section~\ref{sec:pade}; exponent arithmetic only. \\
Pad\'e remainder positivity and decay & Not treated & Section~\ref{sec:pade}. \\
\bottomrule
\end{tabular}
\end{table}

\paragraph{Reading map.}
Sections~\ref{sec:corridor} and~\ref{sec:tail} are the content;
Sections~\ref{sec:sevenhalves} and~\ref{sec:pade} record what is and is not
formalised of two external routes.  Section~\ref{sec:open} states the
remaining obligations.  Linked phrases open the corresponding Lean
declaration at the pinned source revision~\commitshort.

\medskip
\noindent\textbf{Keywords.} irrationality; Lambert series; rational base;
Pad\'e approximation; Lean~4.
\qquad
\textbf{MSC 2020.} 11J72 (primary); 11J82, 68V20 (secondary).

\section{The coordinatewise corridor}\label{sec:corridor}

We isolate the arithmetic that the clearing scheme leaves behind, in a form
that does not mention the series.

\begin{definition}[coordinatewise corridor]\label{def:corridor}
Let $a,b,N,K,Q,D$ be natural numbers.  Say that $(a,b,N,K,Q,D)$ is a
\emph{coordinatewise corridor} when
\[
 a>0,\qquad Q>0,\qquad D>0,\qquad D\le N+K,\qquad
 a^{K}\mid QD,\qquad Q\,b^{\,N+K+1}<a^{\,K+1}.
\]
\end{definition}

\noindent The reading is: $a$ and $b$ are the numerator and denominator of the
base, $N$ is the shift, $K$ is the width of the cleared window, $Q$ is the
accumulated clearing factor, and $D$ is the final coefficient being cleared.
The bound $D\le N+K$ is the only property of the coefficient used; for the
divisor-counting coefficient it holds because $d(n)\le n$.  The divisibility
$a^{K}\mid QD$ is the requirement that clearing succeeded coordinatewise, and
the last inequality is the tail estimate that makes the trapped integer smaller
than~$1$.  The definition is the
\lword{Erdos1049/RationalBaseLambert.lean}{41}{CoordinatewiseCorridor}{corridor
predicate}.

\begin{theorem}[corridor bound]\label{res:corridorbound}
If $(a,b,N,K,Q,D)$ is a coordinatewise corridor, then
\[
 b^{\,N+K+1}<a\,(N+K).
\]
\end{theorem}

\begin{proof}
Since $QD>0$ and $a^{K}\mid QD$, we have $a^{K}\le QD$, and $D\le N+K$ gives
$a^{K}\le Q(N+K)$.  Hence
\[
 Q\,b^{\,N+K+1}<a^{\,K+1}=a^{K}\cdot a\le Q(N+K)\cdot a,
\]
and cancelling the positive factor $Q$ gives the claim.
\end{proof}

\noindent Formalised as the
\lword{Erdos1049/RationalBaseLambert.lean}{49}{coordinatewiseCorridor_implies_pow_lt_linear}{power-versus-linear
consequence}.

The inequality of Theorem~\ref{res:corridorbound} is where the integer and
rational cases part.  At $b=1$ it reads $1<a(N+K)$, which holds for every
$a\ge2$ and every nonempty window; the corridor imposes no obstruction at all.
At $b\ge2$ the left side is exponential in $N+K$ and the right side is linear,
so the corridor can survive only for small $N+K$.  At $b=2$ the crossing has
already happened at the smallest admissible window.

\begin{proposition}\label{res:exp}
For every natural number $x\ge2$ we have $3x<2^{\,x+1}$.
\end{proposition}

\begin{proof}
Induction from $x=2$, where $6<8$.  For the step, $2^{x+1}\ge2^{2}>3$ when
$x\ge1$, so $3(x+1)=3x+3<2^{x+1}+2^{x+1}=2^{x+2}$.
\end{proof}

\noindent Formalised as the
\lword{Erdos1049/RationalBaseLambert.lean}{70}{three_mul_lt_two_pow_succ}{exponential
comparison}.

\begin{theorem}[no corridor at base $3/2$]\label{res:nocorridor}
For all $N\ge1$ and $K\ge1$ and all natural $Q,D$, the tuple
$(3,2,N,K,Q,D)$ is not a coordinatewise corridor.
\end{theorem}

\begin{proof}
A corridor would give $2^{\,N+K+1}<3(N+K)$ by
Theorem~\ref{res:corridorbound}, contradicting Proposition~\ref{res:exp}
applied to $x=N+K\ge2$.
\end{proof}

\noindent Formalised as the
\lword{Erdos1049/RationalBaseLambert.lean}{83}{threeHalves_no_coordinatewiseCorridor}{corridor
exclusion at three halves}.

\paragraph{What this section proves, and what it does not.}
Theorem~\ref{res:nocorridor} is a statement about Definition~\ref{def:corridor}
and about nothing else.  It says that the arithmetic left over by a literal
coordinatewise clearing at base $3/2$ is inconsistent, and therefore that the
integer-base argument does not transfer in that form.  It does not bound the
denominator of $F(3/2)$, it does not show that $F(3/2)$ is irrational, and it
does not show that $F(3/2)$ is rational.  It also does not rule out a clearing
scheme of a different shape: the corridor fixes one divisibility pattern and
one tail inequality, and an argument organised differently is not covered by
the hypothesis.

\section{The cleared tail and the denominator tax}\label{sec:tail}

Theorem~\ref{res:nocorridor} says that one scheme fails.  This section
identifies the quantity responsible, as an exact recurrence.

Throughout, $r,s,B,F$ are rationals with $r\ne0$, and $c:\N\to\Q$ is arbitrary.
Define the prefix and the cleared tail state by
\[
 P_N=\sum_{m=0}^{N-1}\frac{c(m+1)\,s^{\,m+1}}{r^{\,m+1}},
 \qquad
 U_N=B\,r^{N}\bigl(F-P_N\bigr).
\]
Thus $P_N$ is the partial sum of $\sum_{n\ge1}c(n)(s/r)^{n}$ through level $N$,
and $U_N$ is the tail of a putative value $F$ after that level, scaled by
$Br^{N}$.  These are the
\lword{Erdos1049/RationalBaseLambert.lean}{96}{rationalBasePrefixQ}{rational-base
prefix} and the
\lword{Erdos1049/RationalBaseLambert.lean}{109}{rationalBaseClearedTailQ}{cleared
tail state}.

\begin{theorem}[cleared-tail recurrence]\label{res:tailrec}
For every $N$,
\[
 U_{N+1}=r\,U_N-B\,c(N+1)\,s^{\,N+1}.
\]
\end{theorem}

\begin{proof}
Expanding $P_{N+1}=P_N+c(N+1)s^{N+1}/r^{N+1}$ and $r^{N+1}=r^{N}\cdot r$ in the
definition of $U_{N+1}$ and clearing the denominator $r^{N+1}$, which is
nonzero, gives the identity.
\end{proof}

\noindent Formalised as the
\lword{Erdos1049/RationalBaseLambert.lean}{115}{rationalBaseClearedTailQ_succ}{cleared-tail
recurrence}.

The recurrence has a linear part $rU_N$ and a forcing term $Bc(N+1)s^{N+1}$.
The integer-base argument works by keeping the state in a bounded window, and
the forcing term is what it must absorb at each step.  Its size is exactly the
point of departure.

\begin{theorem}[the forcing term]\label{res:forcing}
Let $s,B$ be natural numbers and $c:\N\to\N$, and put
$G_N=B\,c(N+1)\,s^{\,N+1}$.
\begin{enumerate}
\item If $s\ge2$, $B\ge1$ and $c(N+1)\ge1$, then $2^{\,N+1}\le G_N$.
\item If $s=1$, then $G_N=B\,c(N+1)$.
\end{enumerate}
\end{theorem}

\begin{proof}
For the first part, $2^{N+1}\le s^{N+1}=1\cdot s^{N+1}\le Bc(N+1)s^{N+1}$,
using $B\,c(N+1)\ge1$.  The second part is the definition with $s=1$.
\end{proof}

\noindent Formalised as the
\lword{Erdos1049/RationalBaseLambert.lean}{132}{twoPow_le_rationalBaseForcingNat}{exponential
lower bound} and the
\lword{Erdos1049/RationalBaseLambert.lean}{146}{rationalBaseForcingNat_one}{integer-base
collapse}.

\paragraph{Reading the two parts together.}
At $s=1$ the forcing term is $Bc(N+1)$, so it grows only as fast as the
coefficient; for the divisor function that is $O(N^{\varepsilon})$ for every
$\varepsilon>0$, and a bounded-state argument has room.  At $s\ge2$ the same
term is at least $2^{N+1}$ whenever the coefficient is nonzero.  This is an
exact lower bound on one quantity, and it is all that is proved.  It is not a
proof that no bounded-state argument exists at $s\ge2$; the theorem that one
particular scheme fails is Theorem~\ref{res:nocorridor}, and the bound here
records the size of the term that scheme would have to absorb.

\section{The height criterion at $7/2$}\label{sec:sevenhalves}

Bundschuh and V\"a\"an\"anen~\cite{bv1994} prove an irrationality criterion for
a family of rational bases cut out by a height condition.  The criterion
applies at $\beta=7/2$.  It is an analytic theorem and it is not formalised
here.  Two of its elementary inputs are.

\begin{proposition}\label{res:sevenhalves}
$2^{18}<7^{7}$, and $\tfrac{7}{18}=\tfrac12-\tfrac19$ in $\Q$.
\end{proposition}

\noindent These are the
\lword{Erdos1049/RationalBaseLambert.lean}{29}{sevenHalves_power_certificate}{integer
power comparison}, namely $262\,144<823\,543$, and the
\lword{Erdos1049/RationalBaseLambert.lean}{35}{sevenHalves_rational_margin}{exact
rational boundary}.

The first is the arithmetic content of the estimate
$\log2/\log7<7/18$; the passage from the integer inequality to the logarithmic
one uses monotonicity of the logarithm, which is not part of the formalised
statement.  The second is the exact form of a margin whose other input is the
strict bound $1/\pi^{2}<1/9$, that is $3<\pi$, which is also not formalised
here.  Three inputs are therefore external: the logarithmic step, the bound on
$\pi$, and the analytic theorem of~\cite{bv1994} itself.  We record the two
exact arithmetic facts separately from the theorem precisely so that the two
cannot be confused; a reader should not take Proposition~\ref{res:sevenhalves}
as a formalisation of the published criterion, because it is not one.

The criterion does not cover $3/2$, and the results of
Sections~\ref{sec:corridor} and~\ref{sec:tail} give no evidence either way
about whether $F(3/2)$ is irrational.

\section{Pad\'e denominator exponents}\label{sec:pade}

A second external route replaces the clearing argument by explicit rational
approximation.  One constructs integer linear forms in $1$ and $F(\beta)$ whose
analytic decay outruns the height of their common denominator.  For a
homogenised construction over integer parameters the proposed common
denominator exponent is $E_n=(3n^{2}-n)/2$.  Since only doubled exponents
occur below, every statement lives over $\Z$ and no parity bookkeeping is
needed; we write $\widetilde{E}_n=2E_n=3n^{2}-n$.

\begin{proposition}[summand bound and exact gap]\label{res:pade}
Put
\[
 \widetilde{P}(n,k)=2\bigl(k(n-k)+nk\bigr)+k(k-1),
\]
\[
 \widetilde{Q}(n,m)=2(n^{2}-n)+j^{2}+2jm+j-m^{2}+3m,
 \qquad j=n-m-1 .
\]
Then, for integers $n,k,m$:
\begin{enumerate}
\item if $0\le k\le n$, then $\widetilde{P}(n,k)\le\widetilde{E}_n$, and the
gap factors as
$\widetilde{E}_n-\widetilde{P}(n,k)=(n-k)(3n-k-1)$;
\item $\widetilde{E}_n-\widetilde{Q}(n,m)=2\bigl(n+m(m-1)\bigr)$ identically;
\item if $n\ge0$ and $m\ge1$, then $\widetilde{Q}(n,m)\le\widetilde{E}_n$.
\end{enumerate}
\end{proposition}

\noindent Part~(1) is the
\lword{Erdos1049/RationalPadeArithmetic.lean}{29}{rationalPadePSummandDenExpTwice_le}{summand
exponent bound}, part~(2) the
\lword{Erdos1049/RationalPadeArithmetic.lean}{51}{rationalPadeQMaxDenExpTwice_gap}{exact
gap identity}, and part~(3) the
\lword{Erdos1049/RationalPadeArithmetic.lean}{59}{rationalPadeQMaxDenExpTwice_le}{maximal
exponent bound}.  The gap in~(2) is an identity in $\Z[n,m]$, so~(3) follows
from $m(m-1)\ge0$ and $n\ge0$; the hypothesis $m\ge1$ names the range in which
the corresponding summand does not vanish, and it is not the sharpest
hypothesis under which the inequality holds.

These are inequalities between polynomials in the exponents.  They establish
that the proposed common denominator is large enough to clear the construction,
and nothing further.  Positivity of the remainder, its rate of decay, and the
comparison of that rate against the denominator height are the analytic
obligations, and none of them is treated here.  A reconstructed identity is not
an irrationality measure until those are closed.

\section{Complements and further questions}\label{sec:open}

Problem~\#1049 is open.  Two obligations are separated by the results above.

\begin{problem}[the base $3/2$]
Construct integer linear forms in $1$ and $F(3/2)$ whose analytic decay beats
the height of their common denominator.  Section~\ref{sec:pade} supplies the
exponent arithmetic for a homogenised construction; the remainder estimate is
untouched.  By Theorem~\ref{res:nocorridor}, the coordinatewise clearing route
is unavailable, so an argument here must be of a different shape.
\end{problem}

\begin{problem}[the published region]
Formalise the elementary parameter inequalities of the criterion
of~\cite{bv1994}, keeping the analytic theorem external and cited.
Proposition~\ref{res:sevenhalves} formalises two of its inputs; the
logarithmic monotonicity step and the bound $3<\pi$ remain.
\end{problem}

The unrestricted question, which rational bases give an irrational value,
remains open, and is not reduced to either problem above.  What the two
sections of formal content change is the status of one route, not the status of
the problem: the integer-base clearing argument fails at $3/2$ for an exact and
elementary reason, and the term responsible is the factor $s^{N+1}$ in
Theorem~\ref{res:tailrec}.

\pdfbookmark[1]{Statements and declarations}{statements-declarations}
\subsection*{Statements and declarations}

\paragraph{Artefact and data availability.} The
\href{\sourceurl}{pinned formal-source revision} contains the Lean sources, the
fixed toolchain, and the library manifest used in the verification.  This
manuscript provides navigation rather than proof authority.

\paragraph{Declaration of generative AI use.} Large-language-model agents were
used throughout development to draft and revise prose, formal proofs, and
software.  The author set the objectives and acceptance criteria, selected and
reviewed the claims, and approved the published version.  The author assumes
responsibility for the accuracy, interpretation, and presentation of the work.
Generative systems are production tools; they are not authors and supply no
independent authority.  Lean checks each proof term against the fixed library
version, and the sources linked here contain no proof placeholders and no
project-defined axioms; Lean does not authorise the exposition, the citation
choices, or the interpretation, for which the author remains responsible.

\paragraph{Funding and competing interests.} This work received no external
funding.  The author declares no competing interests.

\paragraph{Acknowledgements.} The problem numbering and status follow the
Erd\H{o}s Problems catalogue maintained by Thomas Bloom~\cite{erdosproblems}.

\appendix
\section{Guide to the formal sources}\label{app:index}

Each linked phrase opens its Lean declaration at the pinned source
revision~\commitshort.  The declarations of this note live in two modules:
the corridor, the cleared-tail recurrence, and the two arithmetic certificates
in the first; the Pad\'e exponent arithmetic in the second.  The link
coordinates are validated against that pinned revision, so they remain correct
as later work moves lines in the working tree.

% The exhaustive source manifest is retained in the authored TeX for automated
% validation, and kept outside the printed note.
\iffalse

\begin{tabular}{@{}ll@{}}
informal statement & linked Lean source\\
corridor predicate & \lrefx{Erdos1049/RationalBaseLambert.lean}{41}{CoordinatewiseCorridor}\\
corridor bound & \lrefx{Erdos1049/RationalBaseLambert.lean}{49}{coordinatewiseCorridor_implies_pow_lt_linear}\\
exponential comparison & \lrefx{Erdos1049/RationalBaseLambert.lean}{70}{three_mul_lt_two_pow_succ}\\
no corridor at $3/2$ & \lrefx{Erdos1049/RationalBaseLambert.lean}{83}{threeHalves_no_coordinatewiseCorridor}\\
rational-base prefix & \lrefx{Erdos1049/RationalBaseLambert.lean}{96}{rationalBasePrefixQ}\\
prefix step & \lrefx{Erdos1049/RationalBaseLambert.lean}{101}{rationalBasePrefixQ_succ}\\
cleared tail state & \lrefx{Erdos1049/RationalBaseLambert.lean}{109}{rationalBaseClearedTailQ}\\
cleared-tail recurrence & \lrefx{Erdos1049/RationalBaseLambert.lean}{115}{rationalBaseClearedTailQ_succ}\\
forcing magnitude & \lrefx{Erdos1049/RationalBaseLambert.lean}{126}{rationalBaseForcingNat}\\
exponential lower bound & \lrefx{Erdos1049/RationalBaseLambert.lean}{132}{twoPow_le_rationalBaseForcingNat}\\
integer-base collapse & \lrefx{Erdos1049/RationalBaseLambert.lean}{146}{rationalBaseForcingNat_one}\\
integer power comparison & \lrefx{Erdos1049/RationalBaseLambert.lean}{29}{sevenHalves_power_certificate}\\
exact rational boundary & \lrefx{Erdos1049/RationalBaseLambert.lean}{35}{sevenHalves_rational_margin}\\
doubled common exponent & \lrefx{Erdos1049/RationalPadeArithmetic.lean}{18}{rationalPadeDenExpTwice}\\
doubled summand exponent & \lrefx{Erdos1049/RationalPadeArithmetic.lean}{23}{rationalPadePSummandDenExpTwice}\\
summand exponent bound & \lrefx{Erdos1049/RationalPadeArithmetic.lean}{29}{rationalPadePSummandDenExpTwice_le}\\
doubled maximal exponent & \lrefx{Erdos1049/RationalPadeArithmetic.lean}{45}{rationalPadeQMaxDenExpTwice}\\
exact gap identity & \lrefx{Erdos1049/RationalPadeArithmetic.lean}{51}{rationalPadeQMaxDenExpTwice_gap}\\
maximal exponent bound & \lrefx{Erdos1049/RationalPadeArithmetic.lean}{59}{rationalPadeQMaxDenExpTwice_le}\\
\end{tabular}
\fi

\begin{thebibliography}{9}
\small
\raggedright
\setlength{\itemsep}{0pt}
\bibitem{erdos1948} P.~Erd\H{o}s,
  \href{https://users.renyi.hu/~p_erdos/1948-04.pdf}{\emph{On arithmetical
  properties of Lambert series}},
  J. Indian Math. Soc. (N.S.) \textbf{12} (1948), 63--66.
\bibitem{erdos1988} P.~Erd\H{o}s, \emph{On the irrationality of certain series:
  problems and results}, in A.~Baker (ed.), \emph{New Advances in Transcendence
  Theory}, Cambridge UP, 1988, pp.~102--109,
  doi:\href{https://doi.org/10.1017/CBO9780511897184.009}{10.1017/CBO9780511897184.009}.
\bibitem{bv1994} P.~Bundschuh and K.~V\"a\"an\"anen,
  \href{https://numdam.org/item/CM_1994__91_2_175_0.pdf}{\emph{Arithmetical
  investigations of a certain infinite product}},
  Compositio Math. \textbf{91} (1994), no.~2, 175--199.
\bibitem{erdosproblems} T.~F. Bloom,
  \href{https://www.erdosproblems.com/1049}{\emph{Erd\H{o}s Problem \#1049}},
  \texttt{erdosproblems.com/1049}, accessed 22 July 2026.
\end{thebibliography}

\end{document}
